%--- iliad ---
% title: Preferences to Rewards
% summary: >-
%   Building from preferences and a minimal set of axioms
%   to a utility function expressible as a sum of discounted rewards:
%   the familiar framing in reinforcement learning.
%--- end ---
\documentclass[11pt]{article}

\usepackage[T1]{fontenc}
\usepackage[utf8]{inputenc}
% Page geometry (compact screen page, 1cm margins) now comes from iliad.sty.
\usepackage{microtype}
\usepackage{mathtools}
\usepackage[round,authoryear]{natbib}

% iliad.sty: local copy first (Overleaf), else the shared ../iliad.sty.
% Provides exercise/solution/definition/theorem/proposition/remark/callout,
% amsmath+amsthm+enumitem+tcolorbox+hyperref+cleveref.
\IfFileExists{iliad.sty}{\usepackage[boxes]{iliad}}{\usepackage[boxes]{../iliad}}

\hypersetup{
  colorlinks=true,
  linkcolor=blue,
  citecolor=blue,
  urlcolor=blue
}

% Axioms are this note's own environment (flat numbering, as in the
% original); everything else comes from iliad.sty.
\newtheorem{axiom}{Axiom}

\newcommand{\Hcal}{\mathcal{H}}
\newcommand{\Ofin}{\mathcal{O}}
\newcommand{\Afin}{\mathcal{A}}
\newcommand{\R}{\mathbb{R}}
\newcommand{\pref}{\succcurlyeq}
\newcommand{\strictpref}{\succ}
\newcommand{\indiff}{\sim}
\newcommand{\Lot}{\Delta}
\newcommand{\emptytraj}{\varepsilon}

\title{Preferences to Rewards}
\author{Fernando E. Rosas}
\date{\today}


\begin{document}

\maketitle

% Fernando's recorded lecture for this day (title auto-queried at build).
\youtube{Zv0EC4-iRjQ}   % axioms of RL

% Fernando's own three, transcribed from the D.1 Doc tab ("idealised agency"
% dump) and confirmed against it by David Q. The bullets these replaced were
% machine-written by the original port -- notes_vNM/main.tex has no outcomes
% list at all -- and were TODO'd out in 460e3d0.
% https://docs.google.com/document/d/1ZI_JIg_-CeMZpZKxlh8TdmVmn6CK3UIbIz_Wl-g9U8o/edit?tab=t.uqpqtto4ldg0
\begin{learningoutcomes}
  \begin{itemize}
    \item Understand the difference between preferences, utility, and reward:
      preferences being a primary, largely uncontroversial notion, and utility
      and rewards being derived notions resting on specific assumptions.
    \item Be able to derive the relationship between various preference
      structures and rationality axioms.
    \item Critically assess alternative notions of rationality, and the
      consequences of dropping various classical decision theory assumptions.
  \end{itemize}
\end{learningoutcomes}

% Overview, not \begin{abstract}: every module opens with \section*{Overview}
% (cf. aixi, solomonoff-induction), and the converter renders it as a heading.
\section*{Overview}
This note develops a short route from preferences over complete trajectories to
expected utility, reward, and discount. We begin with preference relations on
deterministic trajectories and explain how completeness and transitivity yield
an ordinal utility representation. We then show how lotteries, together with
the von Neumann--Morgenstern axioms, produce a cardinal utility over
trajectories, and we clarify the distinction between ordinal preference utility
and vNM utility. Next, following Bowling et al., we add a fifth temporal axiom
that is necessary and sufficient for a recursive representation in terms of
local rewards and discounting. Finally, we explain why reward is not unique:
different reward functions can encode the same utility or the same preference
ordering, affine changes of utility induce corresponding changes of reward, and
potential-based shaping provides a canonical example of reward equivalence.

\section{Introduction}

What does it mean for a system to have a goal? In sequential decision making,
the object of evaluation is often not a single isolated choice or prize but a
whole trajectory: a complete history of states, observations, actions, and
consequences unfolding
through time. Before asking how goals are achieved, it is worth asking a more
basic question: what does it mean to prefer some trajectories over
others, and what follows from that?

Recent work in reinforcement learning has revived this fundamental question,
treating preferences over histories as primitive and asking what additional
assumptions are needed before one can recover scalar reward functions. The
present note provides an introduction to these ideas, which proceeds in three
stages.
\begin{enumerate}[label=(\arabic*)]
    \item First, one asks for a numerical representation of how whole
    trajectories are ranked.
    \item Second, once one allows lotteries over trajectories, one asks when
    these lotteries can be ranked by the expectation of that trajectory
    utility.
    \item Third, one asks what additional requirements are needed in order to
    decompose this expected utility into rewards assigned at each time step.
\end{enumerate}
The first and second steps are closely related to the classical theory developed
by von Neumann and Morgenstern \citep{vonneumann1944theory}. The third asks
what extra temporal structure is needed before utility over whole trajectories
can be decomposed into stagewise rewards, following the line of work developed
in modern reinforcement learning by \citet{pitis2019rethinking},
\citet{shakerinava2022utility}, and \citet{bowling2023settling}.

\section{Preferences over trajectories}

Let $\Ofin$ be a finite set of
observations and $\Afin$ a finite set of actions. A one-step interaction is
then given by $t=(o,a)\in \Ofin\times \Afin$. 
For each $n\in\mathbb{N}_{\geq 0}$, define the space of trajectories of length
$n$ by
\[
\Hcal_n \coloneqq (\Ofin\times \Afin)^n.
\]
We write $\emptytraj$ for the unique trajectory of length $0$. The space of all
\emph{finite} trajectories is
\[
\Hcal^* \coloneqq \bigcup_{n=0}^\infty \Hcal_n.
\]
A typical element of $\Hcal^*$ has the form $h=(o_1,a_1,o_2,a_2,\dots,o_n,a_n)$.
We will keep the notation $\Hcal^*$ for the set of all finite trajectories
throughout.
%For $h\in\Hcal^*$ and $t\in \Ofin\times \Afin$, one may form a longer trajectory
%either by appending or prepending $t$. To stay close to Bowling et al.'s
%notation, we write $t\cdot h$ for the trajectory obtained by prepending $t$ to
%$h$.

\begin{definition}[Preference]
A preference relation on $\Hcal^*$ is a binary relation $\pref$ where
\[
h \pref h'
\]
means that trajectory $h$ is judged at least as good as trajectory $h'$.
\end{definition}

From $\pref$ we derive the usual companion relations:
\[
h \indiff h' \iff h \pref h' \text{ and } h' \pref h,
\qquad
h \strictpref h' \iff h \pref h' \text{ and not } h' \pref h.
\]
We call $\indiff$ `indifference', as an agent has no reason to prefer one over
the other.

At this point, $\pref$ has no properties whatsoever.
One may naturally wonder what kinds of properties it is reasonable to require
of $\pref$, and what follows from them --- which is what we study in the next
sections.


\section{When are preferences problematic?}

Can a preference relation be intrinsically `bad'? The relevant concern here is
whether it leads to some form of self-defeat, avoidable loss, or failure of
coherent behaviour. The discussion in the decision-theory literature suggests at
least three grades of concern.

\begin{enumerate}

\item \textit{Representability failures.}
A first and weakest concern is that a preference relation may fail to be
representable in a convenient form. This does not, by itself, imply that
the preference is irrational. %Incompleteness, for instance, may simply reflect
%genuinely plural or unresolved values; one may refuse to rank some pairs of
%trajectories without making any mistake. Similarly, later in the note we will
%see that violating independence blocks an expected-utility representation, and
%that without the fifth axiom one cannot in general write utility over whole
%trajectories as stagewise reward. These are genuine representational limits, but
Representation failures may merely be inconvenient, but they become more
significant when they are symptoms of deeper issues of the kinds described next
\citep{aumann1962utility,fishburn1970utility}.

\item \textit{Self-defeat and avoidable loss.}
A more serious concern is that preferences may guide choice poorly --- as judged by the
agent's own interest. One important case is \emph{static self-defeat}: choosing an
option that is worse than another available one, or adopting a policy that is
systematically improvable. 
A classic case of suboptimality is dominance: one option or policy dominates
another when it is at least as good in every relevant respect and strictly
better in some, so choosing the dominated option is a clear mistake
\citep{kreps1988notes,mascolell1995micro}. A second case is \emph{diachronic
self-defeat}: a plan that the agent endorses now is predictably
undone later in a way that leaves the agent worse off overall.

\item \textit{Vulnerability.}
The most vivid coherence arguments show that a collection of individually
acceptable choices can be combined into a guaranteed loss. Dutch-book arguments
play this role for credences; money-pump arguments play the analogous role for
preferences. The standard example is a preference cycle
\[
h_1 \strictpref h_2,\qquad h_2 \strictpref h_3,\qquad h_3 \strictpref h_1.
\]
If the agent is willing to pay a small fee to move each time to a strictly
preferred trajectory, then an adversary can guide it around the cycle and back
to where it started, poorer than before \citep{gustafsson2010money}. This is
why intransitivity is usually regarded as a particularly severe pathology: it is
not merely hard to represent, but also vulnerable to exploitation under natural
trading assumptions.

\end{enumerate}

\vspace{3pt}
\begin{callout}

\textbf{Coherence is not selection.}\par\smallskip

It is useful to distinguish three kinds of formal results. A
\emph{representation result} says that if preferences satisfy certain axioms,
then they can be written in a particular mathematical form; the vNM and Savage
theorems being classical examples. A \emph{coherence result} is different: it
links violations of some constraint to a penalty such as a Dutch book, money
pump, dynamic inconsistency, or dominated choice. Complete-class and
admissibility theorems belong more naturally in this second family than in the
first, since they characterize undominated decision rules rather than utility
representations. A \emph{selection result} is different again: it adds a story
about some optimization process --- such as market competition, evolution, or
training dynamics --- and argues that systems lacking a certain property tend to
be selected against. In short, representation concerns \emph{form}, coherence
concerns \emph{vulnerability or domination}, and selection concerns
\emph{survival under pressure}.

Note that representation, coherence, and selection arguments are related, but not identical. 
In general, a coherence argument is a
within-agent claim: it says that if a single agent violates some structural
constraint, then the agent is vulnerable to a penalty such as a money pump,
Dutch book, or dominated choice. A selection argument is different: it asks
whether agents lacking that property would tend to disappear under some broader
optimization pressure, such as market competition, training dynamics, or
evolutionary selection. 
A coherence result may help motivate a selection story, yet it does not by itself
show that realistic environments actually select against the offending
preference pattern. Similarly, a representational failure with no plausible
selection story may still be mathematically interesting while being less central
for explaining the structure of real agents.

For the purposes of this note, the key point is that these notions of badness do
not all coincide. A preference can fail standard representation without being
exploitably incoherent, and not every departure from expected utility is thereby
problematic. Nevertheless, the axioms studied below are useful because --- as we
will see --- they rule out several undesirable properties. 

\end{callout}

\newpage

\section{Completeness, transitivity, and ordinal utility}

Two structural conditions are especially important.

\begin{definition}[Completeness]
A preference relation $\pref$ on $\Hcal^*$ is \emph{complete} if for every pair
$h,h'\in \Hcal^*$,
\[
h\pref h' \quad \text{or} \quad h'\pref h \quad \text{(or both)}.
\]
\end{definition}

\begin{definition}[Transitivity]
A preference relation $\pref$ on $\Hcal^*$ is \emph{transitive} if for every
$h,h',h''\in \Hcal^*$,
\[
h\pref h' \quad \text{and} \quad h'\pref h'' \quad \Longrightarrow \quad h\pref h''.
\]
\end{definition}

Completeness says that the agent can compare any two trajectories. This is a
claim about \emph{comparability}, not about certainty: it says that the
preference relation returns a verdict on every pair, not that the agent knows
everything about the consequences of those trajectories. Transitivity says that
these verdicts fit together consistently across chains of comparison. For
example, if $h_1\pref h_2$ and $h_2\pref h_3$, then transitivity requires
$h_1\pref h_3$ as well.

A preference relation satisfying both conditions is often called a
\emph{weak order} in economics and decision theory
\citep{fishburn1970utility,kreps1988notes,mascolell1995micro}. On a countable
domain such as $\Hcal^*$, weak orders admit an ordinal utility
representation. More general representation theorems on richer domains go back
to the classic work of \citet{debreu1954representation}.

\begin{proposition}[Ordinal representation on the trajectory space]
\label{prop:ordinal_representation}
A preference relation $\pref$ on $\Hcal^*$ is complete and transitive if and
only if there exists a function $u\colon \Hcal^* \to \R$ such that
\[
h \pref h' \iff u(h)\geq u(h')
\qquad
\text{for all } h,h'\in\Hcal^*.
\]
\end{proposition}

\begin{proof}
If such a function $u$ exists, then completeness and transitivity are inherited
from the total order on $\R$. Conversely, if $\pref$ is complete and
transitive, then trajectories can be grouped into indifference classes, where
each class contains all trajectories tied with one another. The quotient
$\Hcal^*/{\indiff}$ of these classes is a countable total order. Any countable
total order can be embedded in $\R$, so we may assign real numbers to the
indifference classes in a way that preserves their order. Composing that
assignment with the quotient map gives the desired utility function on
trajectories.
\end{proof}

This utility is \emph{ordinal}: only the ranking matters. Any strictly
increasing transformation of $u$ represents the same preference relation. So
ordinal utility lets us encode the order of deterministic trajectories, but it
does not yet tell us how to compare risky mixtures of them. The numbers
themselves have no independent meaning beyond the order they induce. If $u$
represents a preference relation, then so does $2u+17$, or $\exp(u)$, provided
the transformation remains strictly increasing. This is why ordinal utility is
best thought of as a numerical \emph{labeling of ranks}, not yet as a measure of
how much one trajectory is preferred to another \citep{fishburn1970utility}.

It is also useful to understand what happens when either condition fails.

\subsection*{What if completeness or transitivity fail}

\paragraph{If completeness fails.}
Then some pairs of trajectories are incomparable. This may be a feature rather
than a bug: the agent may genuinely refuse to rank certain alternatives because
its values are plural, under-specified, or context-sensitive. But once
incomparability is allowed, no single real-valued utility function can exactly
represent the relation in the sense of
\cref{prop:ordinal_representation}, because any two real numbers are
themselves comparable. One must then move to a different formalism, such as
partial orders, sets of utility functions, or multi-criteria representations
\citep{aumann1962utility}.

\paragraph{If transitivity fails.}
Then local pairwise judgments need not assemble into a global ranking. In the
simplest case one obtains a cycle
\[
h_1 \strictpref h_2,\qquad h_2 \strictpref h_3,\qquad h_3 \strictpref h_1.
\]
No scalar utility can represent such a cycle, since it would require
\[
u(h_1)>u(h_2)>u(h_3)>u(h_1),
\]
which is impossible.
The logical problem is already serious: there is no single global ranking of the
options. Under the additional behavioral assumption that the agent is willing to
pay a small fee to move from any option to a strictly preferred one, this
logical failure becomes operational through the \emph{money pump} problem.
Suppose an agent currently has $h_1$ and is willing to pay a small fee
$\varepsilon>0$ each time it moves to a strictly preferred trajectory. The
cycle above licenses the sequence
\[
h_1 \to h_2 \to h_3 \to h_1,
\]
with the agent paying $\varepsilon$ at each step. At the end it is back where
it started, but poorer by $3\varepsilon$. Repeating the cycle pumps away
arbitrarily much money \citep{gustafsson2010money}.

There is also an interesting geometric view on transitivity. This paragraph is
not needed for the rest of the note, so readers meeting these ideas for the
first time can safely treat it as an optional aside. Fix a finite menu
$V=\{h_1,\dots,h_m\}\subset \Hcal^*$ and draw the complete graph whose vertices
are the trajectories in $V$. Encode pairwise comparisons by an antisymmetric
edge flow $X$:
\[
X_{ij} = -X_{ji},
\]
where $X_{ij}>0$ means that $h_i$ is preferred to $h_j$, while $X_{ij}<0$ means
the reverse. If preferences come from a utility function $u$, then each edge
weight is just a utility difference,
\[
X_{ij}=u(h_i)-u(h_j),
\]
so $X$ is a discrete gradient field. In the language of the discrete
Helmholtz--Hodge decomposition, any edge flow can be split into a gradient part,
which comes from a potential, and a cyclic part, which records genuine loops.
On the complete comparison graph there is no extra harmonic remainder, so
inconsistency is entirely captured by the cyclic component
\citep{jiang2011statistical}. A three-cycle corresponds exactly to a nonzero
discrete curl:
\[
X_{ij}+X_{jk}+X_{ki}\neq 0.
\]
Thus transitive preferences are precisely the \emph{potential} part of the
decomposition, with the potential given by the utility, while intransitive
cycles show up as the rotational or cyclic part. If this language feels
abstract, the key takeaway is simple: utility means all local comparisons come
from one global ranking, whereas cycles are the leftover pattern that cannot be
explained by any single scalar potential.

\section{Lotteries and the von Neumann--Morgenstern axioms}
\label{sec:lotteries-vnm}

Up to this point, we have just considered preferences over a countable set
$\Hcal^*$ without considering stochasticity.
To introduce uncertainty, we expand the domain from trajectories to lotteries
over trajectories.

Let $\Lot(X)$ for $X\subset\Hcal^*$ denote the set of probability
distributions over $X$. An element $\mu\in\Lot(X)$ can be expressed
as
\[
\mu = \sum_{i} p_i \delta_{x_i},
\]
which should be read as a lottery that yields outcome $x_i$ with probability $p_i$.
The outcome $x$ is identified with the degenerate lottery $\delta_x$.
In this way, preferences over outcomes can be viewed as a special case of
preferences over lotteries.
Furthermore, for $\mu,\nu \in \Lot(X)$ and $\lambda \in [0,1]$, write
\[
\lambda \mu + (1-\lambda)\nu
\]
for the compound lottery that first flips a coin with bias $\lambda$, then samples from $\mu$ or $\nu$ accordingly.

The von Neumann--Morgenstern (vNM) framework studies weak preference relations
$\pref$ over $\Lot(X)$ and asks when such preferences admit an expected-utility
representation \citep{vonneumann1944theory}. The framework considers four axioms
on preferences over lotteries.

\begin{axiom}[Completeness]
For all $\mu,\nu\in\Lot(X)$, either $\mu\pref \nu$ or $\nu\pref \mu$ (or both).
\end{axiom}

\begin{axiom}[Transitivity]
For all $\mu,\nu,\xi\in\Lot(X)$, if $\mu\pref\nu$ and $\nu\pref\xi$, then $\mu\pref\xi$.
\end{axiom}

\begin{axiom}[Continuity]
For all $\mu,\nu,\xi\in\Lot(X)$ with $\mu\pref\nu\pref\xi$, there exists $\lambda\in[0,1]$ such that
\[
\nu \indiff \lambda \mu + (1-\lambda)\xi.
\]
\end{axiom}

\begin{axiom}[Independence]
For all $\mu,\nu,\xi\in\Lot(X)$ and $\lambda\in(0,1)$,
\[
\mu\pref\nu
\quad\Longleftrightarrow\quad
\lambda \mu + (1-\lambda)\xi \pref \lambda \nu + (1-\lambda)\xi.
\]
\end{axiom}

We already know about completeness and transitivity from the previous section; the
new players are continuity and independence.
Continuity says intermediate prospects admit a break-even mixture between better and worse ones.
Independence says that if $\mu$ is preferred to $\nu$, then mixing both with the same background lottery $\xi$ should not reverse that preference.

The natural question is therefore when a preference over lotteries can be represented by the expectation of a utility function on trajectories:
\[
U(\mu)=\sum_{x\in X}\mu(x)u(x).
\]
This formulation says that the value of a lottery is the probability-weighted
average of the utilities of its possible trajectories.
In its simplest finite-outcome form, this question is answered by the
celebrated vNM theorem.


\begin{theorem}[von Neumann--Morgenstern]
Let $X\subseteq \Hcal^*$ be finite. A weak preference relation $\pref$ on
$\Lot(X)$ satisfies completeness, transitivity, continuity, and independence if
and only if there exists a function $u\colon X\to\R$ such that for all lotteries
$\mu,\nu\in\Lot(X)$,
\[
\mu\pref\nu
\quad\Longleftrightarrow\quad
\sum_{x\in X}\mu(x)u(x)\geq \sum_{x\in X}\nu(x)u(x).
\]
Moreover, $u$ is unique up to positive affine transformations:
\[
u'(x)=a+bu(x), \qquad b>0.
\]
\end{theorem}

\begin{proof}
We first prove the easy direction: if preferences admit an expected-utility
representation, then they satisfy the four axioms.

Completeness and transitivity follow immediately from the total order on
$\R$. Continuity follows because if $\mu\pref\nu\pref\xi$ then 
$U(\mu)\geq U(\nu)\geq U(\xi)$, so one can choose $\lambda\in[0,1]$ such that
\[
U(\nu)=\lambda U(\mu)+(1-\lambda)U(\xi).
\]
Hence
\[
\nu \indiff \lambda \mu +(1-\lambda)\xi.
\]
Finally, independence follows from linearity:
\[
U(\lambda \mu +(1-\lambda)\xi)
=
\lambda U(\mu)+(1-\lambda)U(\xi).
\]
Therefore
\[
\mu\pref\nu
\quad\Longleftrightarrow\quad
\lambda \mu +(1-\lambda)\xi \pref \lambda \nu +(1-\lambda)\xi.
\]

We now prove the converse direction, namely that the four axioms imply an
expected-utility representation.

\begin{enumerate}[label=\textbf{\arabic*.}, leftmargin=*]
\item \emph{Pick best and worst outcomes (completeness and transitivity).}
Because $X$ is finite and the restriction of $\pref$ to degenerate lotteries is
complete and transitive, either all outcomes in $X$ are indifferent or there
exist $b,w\in X$ such that
\[
b\pref x \pref w
\qquad
\text{for all } x\in X.
\]
If all outcomes are indifferent, then the constant utility function already
represents the preference relation. So we may assume $b\strictpref w$.

\item \emph{Calibrate each outcome against $b$ and $w$ (continuity).}
For $x=b$ set $u(x)=1$, and for $x=w$ set $u(x)=0$. For
$b\strictpref x \strictpref w$, continuity gives some $u(x)\in(0,1)$ such that
\[
x \indiff u(x)b+\big(1-u(x)\big)w.
\]
Thus every outcome is indifferent to a lottery over the two benchmark outcomes
$b$ and $w$.

\item \emph{Compare the benchmark lotteries (independence and transitivity).}
We first show that for $x,x'\in X$,
\[
x\pref x'
\quad\Longleftrightarrow\quad
u(x)\geq u(x').
\]

For $\lambda\in[0,1]$, write
\[
C_\lambda \coloneqq \lambda b +(1-\lambda)w.
\]
We claim that if $\alpha>\beta$, then $C_\alpha \strictpref C_\beta$. Indeed,
if $\alpha=1$, then $C_\alpha=b$, and independence applied to
$b\strictpref w$ with mixing weight $\beta$ and background lottery $b$ gives
\[
b \strictpref \beta b +(1-\beta)w = C_\beta.
\]
So it remains to consider the case $\alpha<1$. Define
\[
\rho \coloneqq \frac{\alpha-\beta}{1-\beta}\in(0,1).
\]
Since $b\strictpref w$, independence gives
\[
\rho b +(1-\rho)C_\beta
\strictpref
\rho w +(1-\rho)C_\beta.
\]
But the left-hand side is exactly $C_\alpha$, while the right-hand side is
$C_\beta$. Hence larger values of $\lambda$ yield strictly better benchmark
lotteries. Together with Step 2 and transitivity, this proves the displayed
equivalence above.

\item \emph{Reduce an arbitrary lottery to a benchmark lottery (independence).}
Let
\[
\mu=\sum_i p_i x_i.
\]
By Step 2, each $x_i$ is indifferent to $C_{u(x_i)}$. Replacing each $x_i$ by
the corresponding $C_{u(x_i)}$ inside $\mu$, one at a time, and using
independence together with transitivity, yields
\[
\mu \indiff \sum_i p_i C_{u(x_i)}.
\]
By ordinary probability arithmetic, the compound lottery on the right reduces
to
\[
\sum_i p_i C_{u(x_i)}
\indiff
U(\mu)b+(1-U(\mu))w,
\qquad
U(\mu)\coloneqq \sum_i p_i u(x_i).
\]

\item \emph{Conclude the expected-utility representation.}
Applying Step 3 to the benchmark lotteries from Step 4, we obtain
\[
\mu\pref\nu
\quad\Longleftrightarrow\quad
U(\mu)\geq U(\nu).
\]
This is exactly the desired expected-utility representation.

\end{enumerate}

It remains to prove affine uniqueness, which uses continuity, independence,
and the representation just obtained. Suppose $u'$ is another
expected-utility representation of the same preference relation. Let
\[
a\coloneqq u'(w),
\qquad
c\coloneqq u'(b)-u'(w)>0.
\]
For any $x\in X$, Step 2 gives
\[
x \indiff u(x)b+(1-u(x))w.
\]
Since $u'$ also represents the same preferences, indifference implies equality
of expected $u'$-value:
\[
u'(x)=u(x)u'(b)+(1-u(x))u'(w)=a+cu(x).
\]
So $u'$ is a positive affine transformation of $u$.

\end{proof}

\begin{exercise}[Guided proof of the von Neumann--Morgenstern theorem]
\label{ex:vnm-theorem}
Assume $X\subseteq \Hcal^*$ is finite and that $\pref$ on $\Lot(X)$ satisfies
completeness, transitivity, continuity, and independence.
\begin{enumerate}[label=(\alph*)]
\item Show that there exist trajectories $b,w\in X$ such that
$b\pref h \pref w$ for every $h\in X$.
\item Using continuity, show that for every $h\in X$ there exists
$u(h)\in[0,1]$ such that
\[
h \indiff u(h)b+(1-u(h))w.
\]
\item Use transitivity to show that if
\[
h \indiff u(h)b+(1-u(h))w
\quad\text{and}\quad
h' \indiff u(h')b+(1-u(h'))w,
\]
then
\[
h\pref h' \quad\Longleftrightarrow\quad u(h)\geq u(h').
\]
\item Let
\[
\mu=\sum_i p_i h_i.
\]
Use independence repeatedly to replace each $h_i$ by the equivalent lottery
$u(h_i)b+(1-u(h_i))w$ and show that
\[
\mu \indiff \sum_i p_i\bigl(u(h_i)b+(1-u(h_i))w\bigr).
\]
\item Reduce the compound lottery in part (d) to a simple lottery and prove
that
\[
\mu \indiff U(\mu)b+(1-U(\mu))w,
\qquad
U(\mu)\coloneqq \sum_i p_i u(h_i).
\]
\item Show that
\[
\mu\pref\nu \quad\Longleftrightarrow\quad U(\mu)\geq U(\nu).
\]
This gives the expected-utility representation.
\item Prove affine uniqueness: if $u'$ also represents the same preference
relation in expected-utility form, then $u'=a+bu$ for some $a\in\R$ and $b>0$.
\item Prove the converse direction: if preferences admit an expected-utility
representation, then they satisfy completeness, transitivity, continuity, and
independence.
\end{enumerate}
\end{exercise}


In contrast to our previous result, vNM utility is \emph{cardinal} up to
positive affine transformations, not merely ordinal. A general
monotone transformation would destroy the expectation formula. The independence
axiom forces exactly the amount of structure needed for linear averaging.


\medskip
\begin{callout}
\textbf{Two notions of utility.}

In economics it is standard to distinguish two different notions of utility
\citep{kreps1988notes,mascolell1995micro}.

The first is \emph{ordinal utility}, also known as \emph{preference utility} or
`F1', which represents preferences over certain outcomes. This is the object
obtained in \cref{prop:ordinal_representation}: if preferences over
deterministic trajectories are complete and transitive, then there exists a
function $u_{\mathrm{ord}}$ such that
\[
h\pref h'
\quad\Longleftrightarrow\quad
u_{\mathrm{ord}}(h)\geq u_{\mathrm{ord}}(h').
\]
Only the ranking matters. Any strictly increasing transformation of
$u_{\mathrm{ord}}$ represents the same preferences \citep{fishburn1970utility}.

The second is \emph{von Neumann--Morgenstern utility}, also called `F2'. This is the function
$u_{\mathrm{vNM}}$ that appears inside the expectation operator when
preferences over lotteries satisfy continuity and independence:
\[
\mu\pref\nu
\quad\Longleftrightarrow\quad
\sum_h \mu(h)u_{\mathrm{vNM}}(h)\geq \sum_h \nu(h)u_{\mathrm{vNM}}(h).
\]
Unlike ordinal utility, $u_{\mathrm{vNM}}$ is unique only up to positive affine
transformations. Its numerical differences therefore carry behavioral content:
they determine which mixtures an agent is willing to accept, and so encode the
agent's attitudes toward risk and gambling structure
\citep{vonneumann1944theory,mascolell1995micro}.

\end{callout}
\medskip


\begin{exercise}[Ordinal versus vNM utility]
\label{ex:ordinal-vs-vnm}
Let $x,y,z\in\Hcal^*$ satisfy $x\strictpref y\strictpref z$.
\begin{enumerate}[label=(\alph*)]
\item Give two different ordinal utility functions $u_1,u_2\colon\Hcal^*\to\R$
that represent the same ranking of $x,y,z$.
\item Explain why these two functions are equally good as ordinal
representations.
\item Suppose in addition that
\[
y \indiff \tfrac12 x + \tfrac12 z.
\]
Show that not every strictly increasing transformation of a vNM utility can
preserve this indifference.
\end{enumerate}
\end{exercise}

\subsection*{If continuity or independence fails}

It is useful to separate the roles of the four vNM axioms:
\begin{itemize}
		\item As before, completeness and transitivity give us an ordinal utility on deterministic trajectories. 
		\item Continuity and independence turn that ordinal picture into an \emph{expected} utility theory over uncertain prospects. 
\end{itemize}
If completeness or transitivity fail, as discussed in the previous section, we
lose hope of describing the preference by a single scalar quantity.
If continuity or independence fails, the preference will generally no longer
admit the linear expectation form.


\paragraph{Failure of continuity.}
To build intuition, let $x\succ y\succ z$ and define
\[
\mu_\lambda=\lambda x + (1-\lambda)z.
\]
Then
\[
\mu_0=z,\qquad \mu_1=x,
\]
so as $\lambda$ increases from $0$ to $1$, the lottery moves from the worse
prospect $z$ to the better prospect $x$. Continuity says, informally, that
preferences vary smoothly along this path. In particular, it says that the
intermediate prospect $y$ can be matched by some mixture of the better and worse
prospects:
\[
y \indiff \lambda x + (1-\lambda)z
\]
for some $\lambda\in(0,1)$. Intuitively, the axiom requires no abrupt jumps
in preference as the mixture probabilities vary.


When continuity fails, some distinctions can become \emph{lexicographic}.
For example, we could have a case where $y\succ \mu_0 = z$, but
$\mu_\lambda\succ y$ for all $\lambda>0$.
This means that $x$ has absolute priority over $y$: as soon as $x$ appears with
any nonzero probability, the lottery becomes strictly better than $y$.
One way to interpret this is that some considerations are given lexical
priority over others. For instance, avoiding catastrophe might outrank ordinary
gains so completely that no finite improvement in ordinary reward compensates
for even an arbitrarily small increase in catastrophic risk.
Such preferences need not be contradictory; they are simply too sharp to be
captured by a single real-valued utility whose expectations are taken in the
usual way.

What breaks in this case is not all representations, but the \emph{real-valued}
vNM representation. If continuity is dropped, one is naturally led to
lexicographic or non-Archimedean representations: for example, ordered pairs or
vectors of utilities compared lexicographically, or utilities with infinitesimal
scales \citep{fishburn1971lexicographic}. In those models, the first coordinate
records the highest-priority consideration, and lower coordinates matter only
when higher ones tie. 


\paragraph{Failure of independence.}

Independence says that common lottery components should cancel. If $\mu\pref\nu$, then mixing both with the same background lottery $\xi$ at the same rate should preserve the ranking:
\[
\lambda \mu +(1-\lambda)\xi \pref \lambda \nu +(1-\lambda)\xi.
\]
So the relative ranking of $\mu$ and $\nu$ should depend only on how they differ, not on the common part they share.

When independence fails, the value of a prospect depends on its \emph{context}. The same local substitution can be attractive in one background and unattractive in another. This is exactly what the Allais pattern shows \citep{allais1953comportement}. Let $h_5\succ h_1\succ h_0$ denote trajectories yielding $5$, $1$, and $0$ million, and consider
\[
A = 1\cdot h_1,
\qquad
B = 0.89\, h_1 + 0.10\, h_5 + 0.01\, h_0,
\]
\[
C = 0.11\, h_1 + 0.89\, h_0,
\qquad
D = 0.10\, h_5 + 0.90\, h_0.
\]
Many people prefer $A$ to $B$, but prefer $D$ to $C$. Under independence this is impossible, because $A$ versus $B$ and $C$ versus $D$ differ only by a common consequence. The reversal shows that certainty is treated as psychologically special: replacing a sure outcome by a tiny risk of getting nothing matters more than expected utility allows.

This is the general lesson of independence failure. Probabilities are no longer aggregated linearly against a fixed utility function on trajectories. Common branches cannot be canceled, and the whole shape of the distribution starts to matter. Decision makers may overweight certainty, distort small probabilities, care about disappointment or regret, or evaluate gains and losses relative to a reference point rather than in absolute terms. This is the route taken by prospect theory and related non-expected-utility models \citep{kahneman1979prospect,machina1982expected}.

From the perspective of sequential choice, independence also matters because it allows one to replace a sublottery by an equivalent one without changing the value of the larger plan. If independence fails, the value of a branch may depend on the branches surrounding it, so local and global evaluations need not line up. Hammond's consequentialist argument shows that, together with dynamic consistency and suitable sequential assumptions, one is pushed back toward independence \citep{hammond1988consequentialist}. But that only shows one route to coherent planning. One may instead keep a richer, non-linear evaluation of lotteries and give up the idea that common consequences are always behaviorally irrelevant.

So the two failures have different meanings. Failure of continuity says that some priorities are infinitely sharp, leading naturally to lexicographic or infinitesimal utility scales. Failure of independence says that uncertainty is evaluated holistically rather than by linear averaging, leading to models in which background risk, certainty, or reference dependence affect choice.


\begin{exercise}[Failures of the vNM axioms]
\label{ex:vnm-failures}
Here we will study the consequences of different axioms.
\begin{enumerate}[label=(\alph*)]
\item Construct a simple incomplete preference relation on three trajectories.
Why can it not be represented by a single real-valued ordinal utility?
\item Construct a three-cycle and explain how it gives rise to a money pump.
\item Give an example of a lexicographic preference over three outcomes and show
that it violates continuity.
\item Write down the Allais pattern from \cref{sec:lotteries-vnm} and explain which axiom it
violates.
\end{enumerate}
\end{exercise}

\section{A fifth axiom: reward and discount}

The vNM theorem gives an expected-utility representation over lotteries of whole
trajectories, but it does not yet tell us that utility can be generated
\emph{locally} from stepwise rewards. 

To investigate the possibility of localizing utility over time, let us
denote by $T \coloneqq \Ofin\times \Afin$ the set of one-step transitions. 
For $t\in T$ and $h\in \Hcal^*$, write
$t\cdot h$ for the trajectory obtained by prepending $t$ to $h$, and extend this
operation to lotteries by
\[
t\cdot \mu \coloneqq \sum_i p_i \delta_{t\cdot h_i}
\qquad
\text{when }
\mu=\sum_i p_i \delta_{h_i}.
\]
%
Now, using a vNM utility $u(h)$ one can always define a continuation-dependent
increment as
\[
r(t;h) \coloneqq u(t\cdot h)-u(h).
\]
However, this increment may depend on the whole continuation $h$. What is missing is a
temporal condition ensuring that the effect of prepending a transition depends
only on that transition itself. 
Following~\cite{bowling2023settling}, this can be explored with the following
additional axiom.

\begin{axiom}[Temporal $\gamma$-indifference]
There exists a function $\gamma\colon T\to[0,1]$ such that for all
$\mu,\nu\in\Lot(\Hcal^*)$ and all $t\in T$,
\[
\frac{1}{1+\gamma(t)}(t\cdot \mu)+\frac{\gamma(t)}{1+\gamma(t)}\nu
\indiff
\frac{1}{1+\gamma(t)}(t\cdot \nu)+\frac{\gamma(t)}{1+\gamma(t)}\mu.
\]
\end{axiom}

Under expected utility, the indifference above is equivalent to
\[
u(t\cdot \mu)+\gamma(t)u(\nu)=u(t\cdot \nu)+\gamma(t)u(\mu),
\]
which can be rearranged as
\[
u(t\cdot \mu)-u(t\cdot \nu)=\gamma(t)\bigl(u(\mu)-u(\nu)\bigr).
\]
So the effect of prepending the same transition $t$ is to rescale the utility
difference between two continuations by a factor $\gamma(t)$. In that sense,
$\gamma(t)$ measures how much the future still matters after the step $t$ has
occurred. When $\gamma(t)=1$, future differences are
preserved without discount at that step. When $0<\gamma(t)<1$, the future still
matters but is damped. When $\gamma(t)=0$, once $t$ has occurred the
continuation no longer affects the comparison. Bowling's axiom is closely
related to the transition-dependent discounting discussed by
\citet{white2017unifying} and \citet{pitis2019rethinking}.

Bowling et al.\ show that adding this axiom to the four vNM axioms is necessary
and sufficient for a discounted-reward representation of preferences.

\begin{theorem}[Markov reward representation, after Bowling et al.]
A weak preference relation $\pref$ on $\Lot(\Hcal^*)$ satisfies completeness,
transitivity, continuity, independence, and temporal $\gamma$-indifference if
and only if there exist functions $u\colon \Hcal^*\to\R$, $r\colon T\to\R$, and
$\gamma\colon T\to[0,1]$
such that $u(\emptytraj)=0$,
\[
u(t\cdot h)=r(t)+\gamma(t)u(h),
\]
and, for all lotteries $\mu,\nu\in\Lot(\Hcal^*)$,
\[
\mu\pref\nu
\quad\Longleftrightarrow\quad
\sum_{h\in\Hcal^*}\mu(h)u(h)\geq \sum_{h\in\Hcal^*}\nu(h)u(h).
\]
Moreover, $r$ is unique up to positive scale, and $\gamma$ is the same function
appearing in the fifth axiom \citep{bowling2023settling}.
\end{theorem}

This theorem sharpens the vNM result in exactly the way reinforcement learning
needs. Utility is no longer an arbitrary scalar attached to a complete
trajectory; it is generated recursively from a local reward and a local weight
on the future. Unrolling the recursion for a trajectory
$h=(t_1,t_2,\dots,t_n)$ gives
\[
u(h)=r(t_1)+\gamma(t_1)r(t_2)+\gamma(t_1)\gamma(t_2)r(t_3)+\cdots+
\Bigg(\prod_{i=1}^{n-1}\gamma(t_i)\Bigg)r(t_n).
\]
Thus the fifth axiom is what allows utility over whole trajectories to be
decomposed into local reward and discounting. Two familiar special cases are worth highlighting. 
\begin{itemize}
		
\item If
$\gamma(t)\equiv \gamma$ is constant, then
\[
u(h)=\sum_{k=1}^n \gamma^{k-1} r(t_k),
\]
which is the standard discounted-return objective. 
\item If
$\gamma(t)\equiv 1$, then
\[
u(h)=\sum_{k=1}^n r(t_k),
\]
so utility is simply the additive cumulative reward. 
\end{itemize}

In summary, the step from vNM utility over whole trajectories to RL-style reward
requires one more axiom. That axiom simultaneously identifies both the local
reward signal and the form of discounting.

\begin{remark}[MDPs, POMDPs, and locality]
The locality result of this section is especially natural in fully
observed MDPs, where one often expects reward to depend only on the current
state transition. In the present note, however, the primitive histories are
sequences of observations and actions, and the corresponding local reward has
the form $r\colon \Ofin\times \Afin\to\R$. In partially observed settings this
can be restrictive: a goal may be Markov in the hidden state, in the agent's
belief state, or in some augmented memory state, without being reducible to a
function of the current observation-action pair alone. In that sense, the fifth
axiom should be read as characterizing when preferences admit a reward that is
local in the chosen representation of experience. If the raw observation stream
is too coarse, one may need to enrich the state description before a Markov
reward representation becomes available \citep{bowling2023settling}.
\end{remark}


\begin{callout}
\textbf{Utility is not reward.}

It is helpful to keep the following four different objects apart.

\begin{itemize}

\item \emph{Preference} is the primitive relation $\pref$. It says only which
trajectories or lotteries are weakly preferred to which others.

\item \emph{Preference utility} (`F1') is an ordinal representation of preferences
over deterministic trajectories. Under completeness and transitivity, it is any
function $u_{\mathrm{ord}}$ such that
\[
h\pref h'
\quad\Longleftrightarrow\quad
u_{\mathrm{ord}}(h)\geq u_{\mathrm{ord}}(h').
\]
It is unique only up to strictly increasing transformations
\citep{fishburn1970utility}.

\item \emph{vNM utility} (`F2') is the stronger, cardinal utility that appears when
preferences over lotteries satisfy the vNM axioms. It is the function
$u_{\mathrm{vNM}}$ for which
\[
\mu\pref\nu
\quad\Longleftrightarrow\quad
\sum_h \mu(h)u_{\mathrm{vNM}}(h)\geq \sum_h \nu(h)u_{\mathrm{vNM}}(h).
\]
It is unique only up to positive affine transformations
\citep{vonneumann1944theory}. Thus `F2' refines `F1': it agrees with it on the
ranking of certain trajectories, but adds the extra structure needed to compare
lotteries.

\item \emph{Reward} is not either of these utilities. A reward function $r(t)$ is a
\emph{local} representation introduced only after adding temporal structure. In
the present framework, reward appears when the fifth axiom allows utility to be
written recursively as
\[
u(t\cdot h)=r(t)+\gamma(t)u(h).
\]
So reward is a way of \emph{decomposing} utility over complete trajectories into
stepwise contributions. It is therefore downstream of preference, and even
downstream of vNM utility. This is why different reward functions can encode the
same utility or the same preference ordering, a point we return to in the next
section \citep{bowling2023settling}.
\end{itemize}

\end{callout}


\begin{exercise}[Guided proof of the Markov reward representation theorem]
\label{ex:markov-reward}
Assume the hypotheses of \cref{ex:vnm-theorem} together with temporal
$\gamma$-indifference, and let $u$ be a vNM utility representation over
trajectories.
\begin{enumerate}[label=(\alph*)]
\item Apply expected utility to temporal $\gamma$-indifference and show that for
all lotteries $\mu,\nu\in\Lot(\Hcal^*)$ and all transitions $t\in T$,
\[
u(t\cdot \mu)+\gamma(t)u(\nu)=u(t\cdot \nu)+\gamma(t)u(\mu).
\]
\item Rearrange part (a) to prove that
\[
u(t\cdot \mu)-\gamma(t)u(\mu)
\]
is independent of the choice of $\mu$.
\item Define
\[
r(t)\coloneqq u(t\cdot \mu)-\gamma(t)u(\mu)
\]
for any $\mu$. Deduce that for every deterministic trajectory $h$,
\[
u(t\cdot h)=r(t)+\gamma(t)u(h).
\]
\item Show that this recursion extends to lotteries by linearity:
\[
u(t\cdot \mu)=r(t)+\gamma(t)u(\mu).
\]
\item Prove the converse direction: if $u$, $r$, and $\gamma$ satisfy
\[
u(t\cdot h)=r(t)+\gamma(t)u(h)
\]
for all $t$ and $h$, and if $u$ extends linearly to lotteries, then temporal
$\gamma$-indifference holds.
\item Unroll the recursion along a finite trajectory
$h=(t_1,\dots,t_n)$ and derive
\[
u(h)=r(t_1)+\gamma(t_1)r(t_2)+\gamma(t_1)\gamma(t_2)r(t_3)+\cdots+
\Bigg(\prod_{i=1}^{n-1}\gamma(t_i)\Bigg)r(t_n).
\]
\item Show that when $\gamma(t)\equiv \gamma$ is constant, the previous formula
reduces to standard discounted return, and when $\gamma(t)\equiv 1$ it reduces
to additive cumulative reward.
\item Show that if $u' = a+bu$, then the corresponding reward must be
\[
r'(t)=br(t)+(1-\gamma(t))a.
\]
Interpret this as a source of reward non-uniqueness.
\end{enumerate}
\end{exercise}

\section{Reward equivalence and shaping}

The preceding theorem shows how to pass from utility over lotteries of
trajectories to a local reward representation. But it also raises an important
question: how unique is that reward? The answer has two parts. If one fixes a
particular utility function $u$ and a particular discount function $\gamma$,
then the reward is determined. But if one changes the numerical representative
of the same preference relation, or allows shaping transformations, then many
different rewards can encode the same underlying ordering.

First observe that once $u$ and $\gamma$ are fixed, the reward is fixed as
well. Indeed, if
\[
u(t\cdot h)=r(t)+\gamma(t)u(h)
\]
for all $t\in T$ and $h\in\Hcal^*$, then necessarily
\[
r(t)=u(t\cdot h)-\gamma(t)u(h),
\]
and the right-hand side must be independent of the continuation $h$. So the
non-uniqueness of reward does not come from ambiguity inside a fixed recursive
representation. It comes from the fact that utility itself is not unique as a
numerical object.

\begin{proposition}[Multiple rewards can encode the same preference relation]
Suppose $u\colon \Hcal^*\to\R$, $r\colon T\to\R$, and
$\gamma\colon T\to[0,1]$ satisfy
\[
u(t\cdot h)=r(t)+\gamma(t)u(h)
\]
for all $t\in T$ and $h\in\Hcal^*$, and that preferences over lotteries are
represented by the expected value of $u$. For any $a\in\R$ and $b>0$, define
\[
u'(h)=a+bu(h)
\qquad\text{and}\qquad
r'(t)=br(t)+(1-\gamma(t))a.
\]
Then $u'(t\cdot h)=r'(t)+\gamma(t)u'(h)$ for all $t\in T$ and $h\in\Hcal^*$, and
the expected value of $u'$ represents the same preference relation as the
expected value of $u$.
\end{proposition}

The proposition shows that reward non-uniqueness already appears at the level
of affine reparameterizations of utility: changing the numerical representative
of the same vNM preference relation generally changes the associated reward
function as well. 
This point is especially simple under the normalization $u(\emptytraj)=0$ used
in the previous section. Then the additive degree of freedom disappears, and one
is left only with positive rescalings:
\[
u'(h)=bu(h),
\qquad
r'(t)=br(t).
\]
So in the normalized setting reward is unique only up to choice of units. 

There is also a second, less trivial kind of non-uniqueness, in which one
changes rewards while leaving the induced trajectory ordering unchanged, first
studied by~\citet{ng1999policy}.

\begin{proposition}[Potential-based shaping changes utility only by a boundary term]
Consider a state-based setting with transitions $t=(s,a,s')$ and a potential
function $\Phi$ on states.
\begin{enumerate}[label=(\alph*)]
\item If $\gamma\equiv 1$ and
$r_\Phi(s,a,s') \coloneqq r(s,a,s')+\Phi(s')-\Phi(s)$, 
then 
\[
\sum_{k=1}^n r_\Phi(t_k)
=
\sum_{k=1}^n r(t_k)+\Phi(s_n)-\Phi(s_0)
\]
for any trajectory $\tau=(t_1,\dots,t_n)$ from $s_0$ to $s_n$. 

\item If $\gamma\in(0,1)$ is constant and
$r_\Phi(s,a,s') \coloneqq r(s,a,s')+\gamma \Phi(s')-\Phi(s)$, 
then
\[
\sum_{k=1}^n \gamma^{k-1}r_\Phi(t_k)
=
\sum_{k=1}^n \gamma^{k-1}r(t_k)-\Phi(s_0)+\gamma^n\Phi(s_n).
\]
\end{enumerate}
\end{proposition}

The previous result shows that shaping $r$ into $r_\Phi$ changes utility only by a boundary term. 
In particular, if all
compared trajectories share the same start state and terminal potential, then
the induced preference ordering is unchanged; if the boundary term vanishes on
all admissible trajectories, then even the numerical utility is unchanged.

The general lesson is that the same preferences admit many different reward
functions, so no single one of them is canonical. Some transformations,
such as positive scaling, merely change the numerical units. Others, such as
potential-based shaping, redistribute value along the trajectory while leaving
the overall preference over complete trajectories unchanged. This is why reward
design is often underdetermined: what matters behaviorally is not a raw reward
function in isolation, but the utility and preference structure it induces.

\begin{exercise}[Reward shaping]
\label{ex:reward-shaping}
Assume additive reward with $\gamma\equiv 1$ and define
\[
r_\Phi(s,a,s')=r(s,a,s')+\Phi(s')-\Phi(s).
\]
\begin{enumerate}[label=(\alph*)]
\item Show that along any finite trajectory $(t_1,\dots,t_n)$ the total shaped
reward differs from the original total reward by the boundary term
$\Phi(s_n)-\Phi(s_0)$.
\item Under what condition on the admissible trajectories does this shaping
leave utility exactly unchanged?
\item Under what weaker condition does it leave only the induced preference
ordering unchanged?
\end{enumerate}
\end{exercise}

\section{Conclusion}

The main lesson of this note is that, in sequential settings, the primitive
object is not a one-step reward but a preference over complete trajectories.
From that starting point one can distinguish several layers of structure.
Completeness and transitivity yield an ordinal utility over deterministic
trajectories. Once lotteries are introduced, continuity and independence refine
that ordinal picture into a von Neumann--Morgenstern utility, which agrees with
the ranking of certain trajectories but carries strictly more information
because it calibrates trade-offs between risky prospects.

This separation of layers clarifies both what expected utility achieves and
what it leaves out. It explains why ``utility'' in economics can refer either
to an ordinal representation of certain preferences or to a cardinal object
suitable for evaluating lotteries. It also shows what fails when particular
axioms are dropped: without completeness or transitivity, scalar representation
itself can fail; without continuity or independence, one may still have
meaningful preferences, but no longer the linear expectation form of vNM. In
that sense, expected utility is not the whole of rationality, but a particular
and mathematically powerful strengthening of it.

The final step is to see that even vNM utility is not yet reward. To obtain a
local reward-and-discount representation one needs additional temporal
structure, captured here by Bowling et al.'s fifth axiom. Under that axiom,
utility over whole trajectories admits a recursive decomposition into local
reward and discounting. But even then reward is not unique: affine changes of
utility induce corresponding changes of reward, and potential-based shaping can
redistribute value along a trajectory while preserving the underlying ordering.
Moreover, the locality of reward depends on the chosen representation of
experience, which is natural in fully observed MDPs but can be restrictive in
partially observed settings unless the state is suitably enriched. The reward
hypothesis, understood in this way, is therefore best read not as the claim
that goals are primitively rewards, but as the claim that sufficiently
structured preferences over trajectories can be represented by rewards.

\bibliographystyle{plainnat}
\bibliography{biblo}


% The exercises that used to be collected here now sit inline, each after the
% material it exercises; the solutions are collected in the appendix instead.
% Fernando's lead-in to the old section is parked verbatim below: its second
% sentence describes an ordering that no longer exists, and rewording it would
% be editing his words. Reinstate or drop it -- don't paraphrase it.
% \newpage
% \section*{Exercises}
%
% The following exercises are meant to help students reconstruct the main results
% of the note for themselves. The last two are guided proof exercises for the two
% central representation theorems.

\section*{Open-ended questions}

\begin{exercise}
\label{ex:open-implications}
Do the results have any prescriptive or descriptive implications? What kinds of agents, with what kinds of preferences should we design? By default, what kind of agents should we expect to obtain via typical training processes?
\end{exercise}

\begin{exercise}
\label{ex:open-weaker-axioms}
Let's assume a superintelligence has preferences over trajectories. Are there weaker versions of the axioms that make it more plausible that these preferences are ``safe for us'' compared to preferences that lead to utilities or even rewards?
\end{exercise}

\begin{exercise}
\label{ex:open-examples}
More generally, are there interesting examples of preferences over trajectories that the students could analyse, that do or do not satisfy some of the axioms?
\end{exercise}


% Solutions live here, each bound to its exercise by the label in its optional
% argument -- which is what lets them sit anywhere at all.
% The header is wrapped so it appears only in the with-solutions PDF: the web
% relocates each solution under its own exercise, and the -nosol PDF has none.
\begin{pdfonly}
\begin{solutionsonly}
\clearpage
\appendix
\section{Solutions}
\label{apx:solutions}
\end{solutionsonly}
\end{pdfonly}

\begin{solution}[ex:vnm-theorem]
\begin{enumerate}[label=(\alph*)]
\item Induction on $|X|$. A single trajectory is its own best and worst
element. Given a maximal $b'$ and minimal $w'$ for $X\setminus\{h\}$,
completeness compares $h$ with each of them and transitivity makes the
better of $\{h,b'\}$ maximal and the worse of $\{h,w'\}$ minimal for $X$.
\item If $b\indiff w$ then every $h$ is indifferent to both and any constant
$u$ works. Otherwise $b\strictpref w$. By (a) we have $b\pref h\pref w$, so
continuity applied to the triple $(b,h,w)$ gives directly some
$\lambda\in[0,1]$ with $h\indiff \lambda b+(1-\lambda)w$; set
$u(h)\coloneqq\lambda$. For uniqueness, independence implies the mixtures
$\alpha b+(1-\alpha)w$ are strictly increasing in $\alpha$ when
$b\strictpref w$, so two different weights cannot both be indifferent to
$h$.
\item Suppose $u(h)\geq u(h')$. Mixture monotonicity gives
$u(h)b+(1-u(h))w \pref u(h')b+(1-u(h'))w$, and chaining the two calibrating
indifferences through transitivity yields $h\pref h'$. If $u(h)<u(h')$ the
same argument gives $h'\strictpref h$. Together:
$h\pref h'\iff u(h)\geq u(h')$.
\item Independence states that $\mu\indiff\nu$ implies
$p\mu+(1-p)\lambda \indiff p\nu+(1-p)\lambda$. View $\mu$ as a mixture in
which $h_1$ appears with weight $p_1$ against the rest; replacing $h_1$ by
its calibrated equivalent $u(h_1)b+(1-u(h_1))w$ therefore leaves the whole
lottery indifferent. Doing this for each $h_i$ in turn — finitely many
steps — and chaining with transitivity gives
$\mu\indiff\sum_i p_i\bigl(u(h_i)b+(1-u(h_i))w\bigr)$.
\item The right-hand side is a compound lottery whose only outcomes are $b$
and $w$; it awards $b$ with total probability $\sum_i p_i u(h_i)=U(\mu)$.
Identifying a compound lottery with the simple lottery it induces gives
$\mu\indiff U(\mu)b+(1-U(\mu))w$.
\item By (e), $\mu$ and $\nu$ are indifferent to calibrated $b$--$w$
mixtures with weights $U(\mu)$ and $U(\nu)$; by mixture monotonicity and
transitivity, $\mu\pref\nu\iff U(\mu)\geq U(\nu)$.
\item Apply the $u'$-representation to the calibration of $h$:
\begin{align*}
u'(h)&=U'\bigl(u(h)b+(1-u(h))w\bigr)=u(h)\,u'(b)+(1-u(h))\,u'(w) \\
&=u'(w)+\bigl(u'(b)-u'(w)\bigr)u(h).
\end{align*}
This is the required form $u'=a+b\,u$, with intercept $u'(w)$ and slope
$u'(b)-u'(w)$, the latter strictly positive because $b\strictpref w$ and
$u'$ represents $\pref$. (Note the statement's scalars $a,b$ are unrelated
to the trajectories $b,w$ of part~(a); the letter $b$ does double duty here,
so read $u'(b)$ as the utility of the best trajectory throughout.)
\item Let $\mu\pref\nu\iff \mathbb{E}_\mu[u]\geq\mathbb{E}_\nu[u]$.
Completeness and transitivity are inherited from the total order on $\R$.
Continuity: $p\mapsto \mathbb{E}_{p\mu+(1-p)\lambda}[u]
=p\,\mathbb{E}_\mu[u]+(1-p)\,\mathbb{E}_\lambda[u]$ is continuous, so upper
and lower contour sets in $p$ are closed. Independence: mixing both sides
with $\lambda$ adds the same $(1-p)\mathbb{E}_\lambda[u]$ and scales the
difference by $p>0$, leaving the comparison unchanged.
\end{enumerate}
\end{solution}

\begin{solution}[ex:ordinal-vs-vnm]
\begin{enumerate}[label=(\alph*)]
\item For instance $u_1(x,y,z)=(2,1,0)$ and $u_2(x,y,z)=(100,5,-3)$.
\item An ordinal representation encodes only the ranking: $u$ represents
$\pref$ iff $u(h)\geq u(h')\iff h\pref h'$, and both functions induce
$x\strictpref y\strictpref z$. Composing with any strictly increasing map
preserves exactly this information, so no ordinal criterion can distinguish
$u_1$ from $u_2$.
\item Under expected utility the indifference forces
$u(y)=\tfrac12 u(x)+\tfrac12 u(z)$. Take the vNM utility
$u(x,y,z)=(1,\tfrac12,0)$ and the strictly increasing map $f(s)=s^3$. Then
$(f\circ u)(x,y,z)=(1,\tfrac18,0)$, while the lottery
$\tfrac12 x+\tfrac12 z$ has expected value $\tfrac12$. Since
$\tfrac18\neq\tfrac12$, the relation represented by $f\circ u$ strictly
prefers the lottery to $y$: the indifference is destroyed. Only positive
affine transformations preserve all such midpoint identities, which is the
uniqueness part of the vNM theorem.
\end{enumerate}
\end{solution}

\begin{solution}[ex:vnm-failures]
\begin{enumerate}[label=(\alph*)]
\item Let $x\strictpref y$, $x\strictpref z$, with $y$ and $z$ incomparable
(neither $y\pref z$ nor $z\pref y$). Any $u\colon\{x,y,z\}\to\R$ satisfies
$u(y)\geq u(z)$ or $u(z)\geq u(y)$ because the reals are totally ordered, so
the relation it represents is complete; incomparability cannot be encoded by
a single real-valued function.
\item Take $x\strictpref y\strictpref z\strictpref x$. Suppose you hold $x$
and are willing to pay some small $\varepsilon>0$ to exchange an item for one
you strictly prefer. Since $z\strictpref x$ you pay to swap $x\to z$; since
$y\strictpref z$ you pay to swap $z\to y$; since $x\strictpref y$ you pay to
swap $y\to x$. You are back where you started, $3\varepsilon$ poorer, and the
cycle can be run forever.
\item Give outcomes two attributes and compare lexicographically (the first
attribute decides unless tied): $A=(1,0)$, $B=(0,1)$, $C=(0,0)$, so
$A\strictpref B\strictpref C$. Extend to lotteries by comparing expected
attribute vectors lexicographically. Continuity would require some
$p\in[0,1]$ with $B\indiff pA+(1-p)C$. But that mixture has attribute vector
$(p,0)$: for every $p>0$ it beats $B$ on the first attribute, and for $p=0$
it is $C\strictpref$-below $B$. No mixing weight produces indifference.
\item With $h_5\strictpref h_1\strictpref h_0$ paying $5$, $1$ and $0$
million, the pattern of \cref{sec:lotteries-vnm} is
\[
A = 1\cdot h_1,
\qquad
B = 0.89\,h_1+0.10\,h_5+0.01\,h_0,
\]
\[
C = 0.11\,h_1+0.89\,h_0,
\qquad
D = 0.10\,h_5+0.90\,h_0,
\]
with $A\strictpref B$ but $D\strictpref C$. Both pairs differ only by a
\emph{common consequence}: replacing $0.89$ of $h_1$ by $0.89$ of $h_0$
turns $A$ into $C$ and $B$ into $D$.
Independence says a ranking is unchanged when the same consequence is mixed
into both sides with the same weight, so the pattern violates independence
(completeness, transitivity, and continuity are all consistent with it).
\end{enumerate}
\end{solution}

\begin{solution}[ex:markov-reward]
\begin{enumerate}[label=(\alph*)]
\item Apply $u$ (in expected-utility form) to both sides of the axiom's
indifference and multiply by $1+\gamma(t)$:
\[
u(t\cdot\mu)+\gamma(t)u(\nu)=u(t\cdot\nu)+\gamma(t)u(\mu).
\]
\item Rearranging, $u(t\cdot\mu)-\gamma(t)u(\mu)=u(t\cdot\nu)-\gamma(t)u(\nu)$
for \emph{all} $\mu,\nu$: the quantity does not depend on which lottery is
used to compute it.
\item By (b) the definition of $r(t)$ is unambiguous. Taking $\mu$ to be the
point mass on a deterministic continuation $h$ gives
$u(t\cdot h)=r(t)+\gamma(t)u(h)$.
\item $t\cdot\mu$ is the pushforward of $\mu$ under prepending $t$, and $u$
is linear on lotteries, so
$u(t\cdot\mu)=\mathbb{E}_{h\sim\mu}\bigl[u(t\cdot h)\bigr]
=\mathbb{E}_{h\sim\mu}\bigl[r(t)+\gamma(t)u(h)\bigr]=r(t)+\gamma(t)u(\mu)$.
\item With the recursion and linearity,
$u(t\cdot\mu)+\gamma(t)u(\nu)=r(t)+\gamma(t)\bigl(u(\mu)+u(\nu)\bigr)$,
which is symmetric under $\mu\leftrightarrow\nu$; hence
$u(t\cdot\mu)+\gamma(t)u(\nu)=u(t\cdot\nu)+\gamma(t)u(\mu)$. Dividing by
$1+\gamma(t)$ shows both mixtures in the axiom have equal expected utility,
and since $u$ represents $\pref$, they are indifferent.
\item Normalise $u(\emptytraj)=0$ (allowed by affine freedom). Induction on
$n$: $u(h)=r(t_1)+\gamma(t_1)\,u\bigl((t_2,\dots,t_n)\bigr)$, and expanding
the inner term yields
\[
u(h)=r(t_1)+\gamma(t_1)r(t_2)+\gamma(t_1)\gamma(t_2)r(t_3)+\cdots+
\Bigl(\prod_{i=1}^{n-1}\gamma(t_i)\Bigr)r(t_n).
\]
\item If $\gamma(t)\equiv\gamma$, the product $\prod_{i<k}\gamma(t_i)$ is
$\gamma^{k-1}$ and $u(h)=\sum_k \gamma^{k-1}r(t_k)$: discounted return. If
$\gamma\equiv 1$ every product is $1$ and $u(h)=\sum_k r(t_k)$: additive
cumulative reward.
\item $r'(t)=u'(t\cdot\mu)-\gamma(t)u'(\mu)
=\bigl(a+b\,u(t\cdot\mu)\bigr)-\gamma(t)\bigl(a+b\,u(\mu)\bigr)
=b\,r(t)+(1-\gamma(t))a$. The same preferences thus admit a family of reward
functions: even with $\gamma$ fixed, $r$ is only determined up to a positive
scaling and an additive shift modulated by $1-\gamma(t)$.
\end{enumerate}
\end{solution}

\begin{solution}[ex:reward-shaping]
\begin{enumerate}[label=(\alph*)]
\item Writing $t_i=(s_{i-1},a_i,s_i)$, the shaped total is
\[
\sum_{i=1}^{n} r_\Phi(t_i)
=\sum_{i=1}^{n} r(t_i)+\sum_{i=1}^{n}\bigl(\Phi(s_i)-\Phi(s_{i-1})\bigr)
=\sum_{i=1}^{n} r(t_i)+\Phi(s_n)-\Phi(s_0),
\]
since the middle sum telescopes.
\item Utility is exactly unchanged iff the boundary term vanishes on every
admissible trajectory, i.e.\ $\Phi(s_n)=\Phi(s_0)$ throughout — for example
when all trajectories start in a fixed $s_0$ and $\Phi$ takes the value
$\Phi(s_0)$ on every reachable terminal state.
\item Only the ordering is at stake if the boundary term is the \emph{same
constant} $c$ for all compared trajectories: then every utility shifts by
$c$ and the induced ranking is untouched. A fixed start state together with
$\Phi$ constant on the reachable terminal states suffices, whatever that
constant is.
\end{enumerate}
\end{solution}

\begin{solution}[ex:open-implications]
Discussion notes, not a unique answer. Prescriptively, the money-pump and
dominance arguments say that an agent which cares about a fungible resource
is under pressure toward completeness and transitivity — coherence is an
attractor for agents that can be exploited otherwise. Descriptively, nothing
guarantees that trained systems satisfy any axiom: gradient descent on
episodic objectives can produce context-dependent, intransitive, or
incomparable preferences, especially off-distribution. A reasonable
expectation is approximate coherence where incoherence was penalised during
training, and no guarantee elsewhere. For design, the trade-off runs both
ways: highly coherent agents are more predictable and analysable but are
exactly the ones for which instrumental-convergence arguments bite; agents
with incomplete or unstable preferences may be harder to exploit into
goal-directed resource acquisition, at the cost of being harder to reason
about.
\end{solution}

\begin{solution}[ex:open-weaker-axioms]
Discussion notes. The natural candidates weaken one axiom at a time.
Dropping \emph{completeness} is the most studied: an agent with incomplete
preferences can remain undecided between continuing and being shut down, so
it is not pushed by coherence arguments toward shutdown-resistance; money
pumps also lose force because the agent may simply refuse trades between
incomparable options. Weakening \emph{continuity} permits lexicographic
safety: ``never cross the constraint, then optimise'' cannot be represented
by a single real-valued utility, which is arguably a feature. Weakening
\emph{independence} allows certainty-favouring (Allais-like) preferences,
which dampen gambling-for-resources behaviour. Weakening \emph{temporal
$\gamma$-indifference} removes the local reward representation: goals about
the shape of a trajectory as a whole (diversity, ``never do X'') stop being
expressible as accumulated reward. The common pattern: each axiom kept buys
a representation theorem and optimisation pressure; each axiom dropped
blocks a coherence-based failure mode while making the agent's behaviour
less analysable.
\end{solution}

\begin{solution}[ex:open-examples]
Discussion notes; instructive examples include: (i) lexicographic
safety-first preferences — violate continuity (see \cref{ex:vnm-failures});
(ii) Pareto/multi-objective preferences, undominated but unaggregated —
violate completeness; (iii) hyperbolic discounting — satisfies the vNM
axioms at a single decision time yet violates temporal
$\gamma$-indifference, hence admits utility but no stationary
reward/discount pair, and is dynamically inconsistent; (iv) the certainty
effect (Allais) — violates independence only; (v) satisficing (``anything
above the threshold is equally fine'') — complete and transitive, so ordinal
utility exists, but indifference plateaus interact oddly with lotteries near
the threshold; (vi) trajectory-shape goals such as ``visit as many distinct
states as possible'' — can satisfy all four vNM axioms (a utility exists)
while violating the temporal axiom, illustrating exactly the gap between
utility and reward.
\end{solution}

\end{document}
